\documentclass[12pt,a4paper]{article}

% Packages
\usepackage[margin=1in]{geometry}
\usepackage{amsmath}
\usepackage{graphicx}
\usepackage{booktabs}
\usepackage{array}
\usepackage{xcolor}
\usepackage{listings}
\usepackage{hyperref}
\usepackage{titlesec}
\usepackage{fancyhdr}
\usepackage{setspace}
\usepackage{caption}
\usepackage{float}

% Colors
\definecolor{codegreen}{rgb}{0,0.6,0}
\definecolor{codegray}{rgb}{0.5,0.5,0.5}
\definecolor{codepurple}{rgb}{0.58,0,0.82}
\definecolor{backcolour}{rgb}{0.95,0.95,0.92}
\definecolor{speedupgreen}{rgb}{0.56,0.93,0.56}
\definecolor{speedupblue}{rgb}{0.68,0.85,0.90}
\definecolor{seqorange}{rgb}{1.0,0.92,0.80}

% Code listing style
\lstdefinestyle{cppstyle}{
    backgroundcolor=\color{backcolour},
    commentstyle=\color{codegreen},
    keywordstyle=\color{blue}\bfseries,
    numberstyle=\tiny\color{codegray},
    stringstyle=\color{codepurple},
    basicstyle=\ttfamily\footnotesize,
    breakatwhitespace=false,
    breaklines=true,
    captionpos=b,
    keepspaces=true,
    numbers=left,
    numbersep=5pt,
    showspaces=false,
    showstringspaces=false,
    showtabs=false,
    tabsize=2,
    language=C++
}
\lstset{style=cppstyle}

% Header/Footer
\pagestyle{fancy}
\fancyhf{}
\rhead{UCS645 -- Parallel and Distributed Computing}
\lhead{TIET, Patiala}
\rfoot{Page \thepage}

% Hyperlink setup
\hypersetup{
    colorlinks=true,
    linkcolor=blue,
    urlcolor=cyan,
}

% Title section formatting
\titleformat{\section}{\large\bfseries}{}{0em}{\thesection.\ }
\titleformat{\subsection}{\normalsize\bfseries}{}{0em}{\thesubsection\ }

% ---------------------------------------------------------------
\begin{document}

% ======================== TITLE PAGE ========================
\begin{titlepage}
    \centering
    \vspace*{1cm}

    {\LARGE\bfseries Parallel and Distributed Computing\\[0.3em]
    UCS645\par}

    \vspace{1.5cm}

    {\Large\bfseries High-Performance Grid-Based Environmental\\
    Spread Modeling Using Multithreaded CPU\\
    Parallelism\par}

    \vspace{2cm}

    {\large\bfseries\underline{Submitted By:}\par}
    \vspace{0.5cm}
    {\large
    Anushka Gupta -- 102303358\\
    Apurv -- 102483087\\
    Ishwanpreet Kaur -- 102313053\\
    Vibhav Garg -- 102303338\par}

    \vspace{1.5cm}

    {\large\bfseries\underline{Submitted to:}\par}
    \vspace{0.3cm}
    {\large
    Dr.\ Saif Nalband\\
    Computer Science and Engineering Department, TIET, Patiala\par}

    \vspace{2cm}

    % Institution logo placeholder -- replace with actual logo file if available
    % \includegraphics[width=0.3\textwidth]{tiet_logo.png}\\[0.5em]
    {\Large\bfseries THAPAR INSTITUTE\\
    OF ENGINEERING \& TECHNOLOGY\\[0.2em]}
    {\normalsize (Deemed to be University)\par}

    \vfill
\end{titlepage}

% ======================== ABSTRACT ========================
\section{Abstract}

This report presents the design, implementation, and performance evaluation of a
parallel grid-based environmental spread simulation using C++ and OpenMP. The
simulation models forest fire propagation on two-dimensional grids of varying sizes
($200\times200$, $400\times400$, and $600\times600$), using probabilistic cell-state
update rules across 200 discrete time steps.

Performance was measured by executing both sequential and parallel implementations
and comparing execution time, speedup, and parallel efficiency. Results demonstrate
that OpenMP-based parallelism with \texttt{collapse(2)} and static scheduling achieves
up to $3.47\times$ speedup on a $600\times600$ grid using 8 threads. The findings
confirm that shared-memory parallelism is effective for embarrassingly parallel grid
workloads, with efficiency characteristics that improve with problem scale.

% ======================== INTRODUCTION ========================
\section{Introduction}

Computational modeling of natural phenomena such as wildfire propagation and heat
diffusion requires intensive numerical computation over spatially distributed data.
Grid-based simulation frameworks are commonly used because each cell in the grid
can be updated independently per time step---making them ideal candidates for
parallel execution.

Parallel and Distributed Computing (PDC) provides essential tools for accelerating
such simulations. Shared-memory parallelism, as enabled by the OpenMP API, allows
multiple CPU threads to collaborate on a single workstation without the complexity of
message-passing frameworks like MPI.

This project implements a forest fire spread model on a 2D grid with three cell states:
\texttt{EMPTY} (0), \texttt{TREE} (1), and \texttt{BURNING} (2). A tree catches fire
if any of its four neighbours is burning; a burning cell becomes empty in the next step.
Starting from a single ignition point at the top-centre of the grid, the simulation runs
for 200 steps. Sequential and parallel variants are benchmarked to quantify the
benefits of parallelism.

% ======================== OBJECTIVES ========================
\section{Objectives}

\begin{itemize}
    \item Implement a forest fire spread simulation on a 2D grid in C++ using
          rule-based cellular automaton logic.
    \item Develop a sequential baseline and an OpenMP-parallel version of the
          simulation update kernel.
    \item Benchmark execution time for grid sizes $200\times200$, $400\times400$,
          and $600\times600$ with thread counts 1, 2, 4, and 8.
    \item Compute and compare speedup and parallel efficiency metrics.
    \item Analyse scalability behaviour and identify performance bottlenecks.
\end{itemize}

% ======================== METHODOLOGY ========================
\section{Methodology}

\subsection{Grid Initialisation}

The grid is initialised with a fixed random seed (\texttt{srand(42)}) ensuring
reproducibility. Each cell is assigned \texttt{TREE} with probability 0.7, or
\texttt{EMPTY} otherwise. A single \texttt{BURNING} cell is placed at position
\texttt{[0][cols/2]}---the centre of the top row---to seed the fire.

\subsection{State Update Rules}

At each simulation step, the current grid is read and a next grid is computed. The
update rules applied to each cell $(i,\,j)$ are:

\begin{itemize}
    \item \texttt{BURNING} $\rightarrow$ \texttt{EMPTY} (the cell burns out)
    \item \texttt{TREE} with a \texttt{BURNING} 4-connected neighbour
          $\rightarrow$ \texttt{BURNING} (fire spreads)
    \item \texttt{TREE} with no burning neighbour $\rightarrow$ \texttt{TREE}
          (unchanged)
    \item \texttt{EMPTY} $\rightarrow$ \texttt{EMPTY} (no regrowth)
\end{itemize}

Because all cells in a step are read from the current grid and written to the next,
updates are fully independent and race-condition-free.

\subsection{Sequential Implementation}

The sequential version iterates over all cells using nested \texttt{for}-loops:

\begin{lstlisting}
// (no directive)
for (int i = 0; i < rows; ++i)
    for (int j = 0; j < cols; ++j) { /* update logic */ }
\end{lstlisting}

\subsection{Parallel Implementation}

The parallel version wraps the same loop body with an OpenMP directive:

\begin{lstlisting}
#pragma omp parallel for collapse(2) schedule(static)
for (int i = 0; i < rows; ++i)
    for (int j = 0; j < cols; ++j) { /* same update logic */ }
\end{lstlisting}

\texttt{collapse(2)} flattens the two-dimensional loop into a single iteration space
before distributing work, enabling better load balance. \texttt{schedule(static)} is
optimal here because every cell performs the same number of operations
(constant-time neighbour check), so equal-size chunks incur minimal overhead.

\subsection{Timing}

Each simulation run is timed using \texttt{std::chrono::high\_resolution\_clock}.
The timer starts immediately before the step loop and stops after all 200 steps.
Grid initialisation time is excluded from measurements. For the parallel case, the
number of threads is set with \texttt{omp\_set\_num\_threads()} before each run.

\subsection{Metrics}

\[
  S = \frac{T_{\text{seq}}}{T_{\text{par}}}
\]

where $T_{\text{seq}}$ is the sequential execution time and $T_{\text{par}}$ is the
parallel execution time for a given thread count. Parallel efficiency is defined as:

\[
  E = \frac{S}{p} \times 100\%
\]

where $p$ is the number of threads.

% ======================== IMPLEMENTATION ========================
\section{Implementation Details}

\begin{itemize}
    \item \textbf{Language \& Compiler:} C++17, compiled with
          \texttt{g++ -O2 -fopenmp}.
    \item \textbf{Data Structure:} \texttt{std::vector<std::vector<int>>} ---
          a standard 2D dynamic array. While a flat 1D vector would offer better
          cache locality, the nested vector was chosen for clarity.
    \item \textbf{Random Seed:} \texttt{srand(42)} ensures identical initial grids
          across all runs, isolating the performance variable to thread count alone.
    \item \textbf{Thread Count Sweep:} The outer loop iterates over sizes
          $\{200,\,400,\,600\}$; for each size, the sequential time is recorded once,
          then parallel runs are performed for thread counts $\{1,\,2,\,4,\,8\}$.
    \item \textbf{Execution Environment:} Kali Linux VM (4.08\,GB RAM),
          \texttt{htop}-confirmed CPU usage up to 84.9\% during peak parallel
          execution, confirming thread activity across all cores.
\end{itemize}

% ======================== RESULTS ========================
\section{Experimental Results}

\subsection{Raw Results Table}

Table~\ref{tab:results} presents the complete set of measured execution times and
computed speedups. Sequential baseline values (Mode = Sequential) are shown for
each grid size.

\begin{table}[H]
\centering
\caption{Execution Time and Speedup --- All Grid Sizes and Thread Counts}
\label{tab:results}
\renewcommand{\arraystretch}{1.3}
\begin{tabular}{>{\centering}p{2.2cm} >{\centering}p{1.5cm} >{\centering}p{2.5cm} >{\centering}p{2cm} >{\centering\arraybackslash}p{2cm}}
\toprule
\textbf{Grid Size} & \textbf{Threads} & \textbf{Mode} & \textbf{Time (s)} & \textbf{Speedup} \\
\midrule
$200\times200$ & 1 & Parallel & 0.4816 & 0.9731 \\
\rowcolor{speedupblue}
$200\times200$ & 2 & Parallel & 0.2855 & 1.6415 \\
\rowcolor{speedupblue}
$200\times200$ & 4 & Parallel & 0.1761 & 2.6609 \\
\rowcolor{speedupblue}
$200\times200$ & 8 & Parallel & 0.1863 & 2.5157 \\
\rowcolor{seqorange}
$200\times200$ & 1 & \textbf{Sequential} & 0.4687 & 1.0000 \\
\midrule
$400\times400$ & 1 & Parallel & 2.1822 & 0.9761 \\
\rowcolor{speedupblue}
$400\times400$ & 2 & Parallel & 1.1578 & 1.8397 \\
\rowcolor{speedupgreen}
$400\times400$ & 4 & Parallel & 0.6971 & 3.0554 \\
\rowcolor{speedupgreen}
$400\times400$ & 8 & Parallel & 0.6390 & 3.3335 \\
\rowcolor{seqorange}
$400\times400$ & 1 & \textbf{Sequential} & 2.1299 & 1.0000 \\
\midrule
$600\times600$ & 1 & Parallel & 5.0303 & 0.9748 \\
\rowcolor{speedupblue}
$600\times600$ & 2 & Parallel & 2.6078 & 1.8804 \\
\rowcolor{speedupgreen}
$600\times600$ & 4 & Parallel & 1.5256 & 3.2142 \\
\rowcolor{speedupgreen}
$600\times600$ & 8 & Parallel & 1.4152 & 3.4651 \\
\rowcolor{seqorange}
$600\times600$ & 1 & \textbf{Sequential} & 4.9036 & 1.0000 \\
\bottomrule
\end{tabular}
\\[0.5em]
\footnotesize
\colorbox{speedupgreen}{\phantom{XX}} Speedup $\geq 3.0\times$ \quad
\colorbox{speedupblue}{\phantom{XX}} Speedup $\geq 2.0\times$ \quad
\colorbox{seqorange}{\phantom{XX}} Sequential baseline
\end{table}

\subsection{Execution Time Analysis}

Figure~1 shows execution time decreasing clearly with thread count for all grid
sizes. The most dramatic reduction is visible in the $600\times600$ case, which
drops from 5.03\,s (1 thread) to 1.42\,s (8 threads). However, going from 4 to 8
threads provides diminishing returns --- for $200\times200$, execution time actually
slightly increases from 0.1761\,s to 0.1863\,s at 8 threads, indicating
thread-management overhead exceeds computation savings at this small scale.

\subsection{Speedup Analysis}

Figure~2 compares achieved speedup against the ideal linear curve. All three grid
sizes track closely to ideal speedup at 2 and 4 threads. At 8 threads, the
$600\times600$ grid achieves $3.47\times$ --- the closest to ideal --- while
$200\times200$ falls to $2.52\times$, demonstrating that larger grids amortise
thread-creation overhead more effectively.

Notably, 1-thread parallel runs consistently show speedup slightly below 1.0
(e.g.\ $0.97\times$), reflecting the small overhead of OpenMP's parallel region
setup even when using a single thread. This is expected and correct behaviour.

\subsection{Parallel Efficiency}

Figure~3 reveals that efficiency degrades gracefully with thread count. For 2
threads, efficiency ranges from 82.1\% ($200\times200$) to 94.0\% ($600\times600$),
reflecting that larger grids better amortise OpenMP setup cost. At 8 threads,
efficiency ranges from 31.4\% ($200\times200$) to 43.3\% ($600\times600$). The
sub-linear growth is consistent with Amdahl's Law --- even a small sequential
fraction (e.g., grid initialisation, swap, I/O) limits peak parallel efficiency.

\subsection{Sequential vs.\ Best Parallel}

Figure~4 directly compares the sequential baseline against the 8-thread parallel
result. The time savings are 0.30\,s ($200\times200$), 1.49\,s ($400\times400$),
and 3.49\,s ($600\times600$). For large grids, parallel execution reduces runtime
to less than 30\% of the sequential time.

% ======================== DISCUSSION ========================
\section{Discussion}

\subsection{Observed Behaviour at 8 Threads ($200\times200$)}

The slight slowdown at 8 threads for the smallest grid (0.1761\,s at 4T vs.\
0.1863\,s at 8T) is attributable to thread-management overhead outweighing the
computation savings. With only 40\,000 cells per step, each thread is assigned fewer
than 5\,000 cells --- too little work to justify 8 thread contexts. This is a
well-known effect in parallel computing known as the \emph{granularity problem}.

\subsection{Choice of \texttt{collapse(2)}}

Using \texttt{collapse(2)} flattens the $i \times j$ iteration space from
$\text{rows} \times \text{cols} = 40\text{K}$--$360\text{K}$ iterations into a
single pool distributed among threads. Without \texttt{collapse}, only the outer
(row) loop would be parallelised --- for a 200-row grid with 8 threads, each thread
would handle 25 rows, providing coarser and potentially imbalanced partitioning.
\texttt{collapse(2)} enables finer-grained load distribution.

\subsection{Choice of \texttt{schedule(static)}}

Static scheduling divides the iteration space into equal contiguous chunks before
execution. This is optimal when each iteration has equal cost --- which holds here
because every cell performs the same $O(1)$ neighbour-check regardless of its state.
Dynamic scheduling would add runtime overhead for no benefit.

\subsection{Memory Layout Consideration}

The grid is stored as \texttt{vector<vector<int>>}, meaning each row is a separate
heap allocation. While this causes pointer indirection on column accesses, it did not
introduce measurable cache-miss overhead at these grid scales. For grids larger than
$10{,}000 \times 10{,}000$, refactoring to a flat \texttt{vector<int>} with manual
row indexing would be advisable.

% ======================== SYSTEM-LEVEL OBSERVATIONS ========================
\section{System-Level Observations}

The \texttt{htop} screenshot captured during execution showed:

\begin{itemize}
    \item CPU core 0 running at 84.9\% --- the primary executing thread.
    \item Cores 1--3 at 4--29\%, consistent with OpenMP distributing work across
          available cores.
    \item Memory usage: 815\,M / 4.08\,GB --- well within limits; the largest grid
          ($600\times600 \times 2$ buffers $\times$ 4 bytes) $\approx$ 2.88\,MB.
    \item Uptime 9 minutes at capture, with 83 tasks and 253 threads active.
\end{itemize}

These observations confirm that the parallel regions genuinely engaged multiple CPU
cores, and the results are not artefacts of single-threaded execution with OMP
overhead.

% ======================== CONCLUSION ========================
\section{Conclusion}

This project successfully demonstrates the application of OpenMP shared-memory
parallelism to a grid-based environmental spread simulation. Key conclusions are:

\begin{itemize}
    \item OpenMP with \texttt{collapse(2)} and static scheduling is highly effective
          for embarrassingly parallel grid workloads, requiring minimal code changes
          over the sequential baseline.
    \item Speedup is grid-size dependent: larger grids ($600\times600$, $3.47\times$)
          consistently outperform smaller grids ($200\times200$, $2.52\times$) at the
          same thread count.
    \item Parallel efficiency decreases with thread count but remains meaningful ---
          43\% at 8 threads for the largest grid --- suggesting the current VM has
          limited physical cores, causing HyperThreading contention at high thread
          counts.
    \item The granularity problem is evident for small grids at high thread counts,
          providing a practical lesson in choosing an appropriate level of parallelism
          for a given problem size.
    \item All timings are reproducible (fixed random seed \texttt{srand(42)}), and the
          parallel logic is provably race-condition-free due to read-only access to the
          current grid.
\end{itemize}

% ======================== FUTURE WORK ========================
\section{Future Work}

\begin{itemize}
    \item Refactor the grid to a flat 1D \texttt{vector} to improve cache locality
          and measure the performance delta.
    \item Extend to GPU acceleration using CUDA or OpenCL to leverage thousands of
          parallel cores.
    \item Implement an MPI-based distributed version for multi-node cluster execution.
    \item Add a fire probability parameter $p_{\text{fire}}$ to make the model
          stochastic and study statistical spread patterns.
    \item Profile with Intel VTune or \texttt{perf} to identify memory bandwidth vs.\
          compute bottlenecks.
\end{itemize}

\end{document}